\section{Polymer geometry}\label{app:polymer-geometry}

We give the geometric estimates used above, following
\cite[Appendix~A]{QED3Manuscript}.

Let $L$ be odd, $M$ a power of $L$, and $\ell$ a multiple of $LM$
with $\ell\ge13LM$. The $M$-cubes have centres in
$M\Z^3/\ell\Z^3$. They partition the lattice sites; their closed
continuum versions may meet along faces and edges. A polymer is a
nonempty union of $M$-cubes connected through faces. Denote its number
of cubes by $|X|_M$. A net in the continuum torus is a connected
finite union of images of specified closed straight segments
in $\R^3$, a point being allowed as a segment.
Its length is the sum of the Euclidean lengths of those segments.
Set
\[
 Md_M(X)=\inf\{\operatorname{length}(N):N
             \text{ is a net meeting every closed cube of }X\},
\]
where the infimum also extends over the finite segment representations.
The same definition is used for finite face-connected unions of
$M$-cubes in $\R^3$, with nets in $\R^3$. Distances are sup distances.

\subsection{Size, diameter, and reblocking}

Put
\[
 c_X=54,\qquad \kappa_0=54\log72,\qquad
 K_0=e^{\kappa_0}/36,\qquad c_{\mathrm{rb}}=c_XK_0.
\]
The following estimates are uniform in $M$ and $\ell$:
\begin{align}
 d_M(X)&\le |X|_M-1,&
 |X|_M&\le c_X(d_M(X)+1),\label{eq:polymer-size-tree}\\
 d(p,q)&\le M(d_M(X)+2)
       &&(p,q\text{ in the closed cubes of }X),
                                  \label{eq:polymer-diameter}\\
 \sum_{X\supset\square}e^{-\kappa_0d_M(X)}&\le K_0,&
 \sum_{X\subset Y}e^{-\kappa_0d_M(X)}&\le K_0|Y|_M.
                                  \label{eq:polymer-counting}
\end{align}
Here $Y$ is any union of $M$-cubes. If $\ell/L$ satisfies the same
side conditions stated above, let $\bar X^L$ be the union of $LM$-cubes containing cubes of
$X$. There is a unique polymer $Y$ on the rescaled torus with
$\bar X^L=LY$, and
\begin{align}
 Ld_M(Y)&\le d_M(X),\label{eq:polymer-reblock-tree}\\
 \sum_{X:\bar X^L=LY}e^{-\kappa'd_M(X)}
      &\le c_{\mathrm{rb}}L^3
           e^{-L(\kappa'-\kappa_0-1)d_M(Y)}
                     \quad(\kappa'\ge\kappa_0).
                                  \label{eq:polymer-reblock-sum}
\end{align}

To prove the first inequality of \eqref{eq:polymer-size-tree}, order
the cubes along a spanning tree of their adjacency graph. A net
meeting one cube can be extended to an adjacent cube by a segment of
length at most $M$ perpendicular to their common face. Starting with
one point and repeating this construction gives a net of length at
most $M(|X|_M-1)$.

For the reverse inequality, subdivide the segments of any net $N$ at
their intersections and discard overlapping pieces. Its total length
does not increase. Doubling the edges of the resulting connected
graph gives a closed traversal of length at most
$2\operatorname{length}(N)$. Cut it into at most
$2\operatorname{length}(N)/M+1$ arcs of length at most $M$.
Each arc lifts to a curve in $\R^3$ whose range in each coordinate
has length at most $M$. Such an interval meets at most three closed
cube intervals, so an arc meets at most $27$ closed $M$-cubes.
Every cube of $X$ meets one of these arcs. Taking the infimum over
$N$ gives
\[
 |X|_M\le27(2d_M(X)+1)\le c_X(d_M(X)+1).
\]
For \eqref{eq:polymer-diameter}, take a net of length less than
$Md_M(X)+\delta$ and points $p',q'$ of the net in the respective
cubes of $p,q$. The path in the net from $p'$ to $q'$ has length
at most that of the net. Hence
$d(p,q)\le M+(Md_M(X)+\delta)+M$; let $\delta\downarrow0$.
The same argument holds in $\R^3$.

A polymer with $n$ cubes containing a fixed cube determines a closed
walk of length $2(n-1)$ starting there: double a spanning tree of its
adjacency graph and traverse it. Choose one such walk for each
polymer. This assignment is injective, since the visited cubes recover
the polymer, and there are at most $6^{2(n-1)}=36^{n-1}$ possible
walks. By the size estimate, $d_M(X)\ge n/c_X-1$. Summing over $n$
gives
\[
 \sum_{X\supset\square}e^{-\kappa_0d_M(X)}
       \le e^{\kappa_0}\sum_{n\ge1}36^{n-1}72^{-n}=K_0.
\]
Sum this bound over the cubes of $Y$ to obtain the second inequality
of \eqref{eq:polymer-counting}.

Every $M$-cube belongs to one $LM$-cube, since $L$ is odd, and adjacent
$M$-cubes belong to the same or to adjacent $LM$-cubes. Thus $\bar X^L$
is connected through faces. Projection by $c\mapsto c/L$ gives its
unique rescaled polymer $Y$. A net for $X$ meets every $LM$-cube of
$\bar X^L$; its rescaling has length divided by $L$. This proves
\eqref{eq:polymer-reblock-tree}. For \eqref{eq:polymer-reblock-sum},
split the decay into rates $\kappa_0$ and $\kappa'-\kappa_0$ and use
\eqref{eq:polymer-reblock-tree}. All terms satisfy $X\subset LY$,
which has $L^3|Y|_M$ fine cubes, so
\[
 \sum_{X:\bar X^L=LY}e^{-\kappa_0d_M(X)}
      \le K_0L^3|Y|_M\le c_XK_0L^3e^{d_M(Y)}.
\]
Combining the two factors proves the stated bound.

\subsection{Nonwrapping lifts}\label{subsec:polymer-lifts}

Let $\operatorname{pr}:\R^3\to\R^3/\ell\Z^3$ be the projection.
\begin{proposition}\label{prop:nonwrapping-polymer-lift}
If $X$ is a polymer and
\begin{equation}\label{eq:nonwrapping-polymer-condition}
 2M(d_M(X)+2)<\ell,
\end{equation}
then $X=\operatorname{pr}(\widetilde X)$ for a finite face-connected
union of $M$-cubes in $\R^3$. The projection is injective on
$\widetilde X$, preserves adjacency of its cubes in both directions,
and $d_M(\widetilde X)=d_M(X)$. The lift is unique up to a translation
in $\ell\Z^3$.
\end{proposition}

\begin{proof}
Put $D=M(d_M(X)+2)$, fix $p_0$ in a closed cube of $X$, and fix its
lift $\widetilde p_0$. The projection is injective on the closed sup
ball $\widetilde B=\{|p-\widetilde p_0|_\infty\le D\}$, since two
points there differ in every coordinate by at most $2D<\ell$.
It maps this ball onto the torus ball $B=\{p:d(p,p_0)\le D\}$.
Its inverse $\varphi:B\to\widetilde B$ is continuous, and is locally
isometric. The diameter estimate places every closed cube of $X$ in
$B$. A segment or a closed $M$-cube lying in $B$ lifts under
$\varphi$ to a segment of the same length or a closed $M$-cube:
the translates $\widetilde B+\ell n$ are pairwise disjoint, and a
connected lift lies in one of them. The same inverse maps shared
faces to shared faces. Thus $\widetilde X=\varphi(X)$ has the
asserted injectivity, adjacency, and connectedness.

Any other face-connected injective lift can be translated so that it
contains the chosen lift of the cube of $p_0$. There are no two
different neighbours of one lifted cube with the same projection:
their centres differ by at most $2M<\ell$. Propagate this observation
along a chain of adjacent cubes. Every cube of the other lift equals
the corresponding cube of $\widetilde X$, proving uniqueness.

Projecting a net for $\widetilde X$ gives a net for $X$ with no greater
length, so $d_M(X)\le d_M(\widetilde X)$. Conversely, let $N$ be a
net for $X$ of length at most $Md_M(X)+M$, and choose its point $q_0$
in the cube of $p_0$. Every point of $N$ is joined to $q_0$ by a
curve of length at most that of $N$, and $d(q_0,p_0)\le M$.
Therefore $N\subset B$. Its lift $\varphi(N)$ has the same length
and meets every cube of $\widetilde X$. Taking an infimizing sequence
of these nets proves the reverse inequality.
\end{proof}

